\documentclass[../main.tex]{subfiles}

\graphicspath{{../image/}{../figure/}}

\begin{document}

\chapter{一维空间中的双曲平衡律方程组}\label{chapter:ch07-一维双曲平衡律组}

第 VII--XVI 章将致力于研究一维空间中的平衡律方程组。研究范围的这种收窄主要是出于必要：就目前而言，多维方程组的理论仍是一片未知之地，充满着引人入胜的问题。无论如何，读者应当记住，某些具有特殊对称性的多维现象，例如波的聚焦，可以在一维空间理论的框架内加以研究。我们将在第 XVII 章回到多维空间。


本章将介绍作为一维空间双曲平衡律方程组理论之基础的许多概念：严格双曲性；Riemann 不变量及其与熵的关系；简单波；真正非线性性及其在古典解破裂中的作用。

为了搭好舞台，本章开篇将给出若干产生于物理学或科学技术其他分支中的一维双曲平衡律方程组的说明性例子。

\subfile{Chapter-07/section-01.tex}
\subfile{Chapter-07/section-02.tex}
\subfile{Chapter-07/section-03.tex}
\subfile{Chapter-07/section-04.tex}
\subfile{Chapter-07/section-05.tex}
\subfile{Chapter-07/section-06.tex}
\subfile{Chapter-07/section-07.tex}
\subfile{Chapter-07/section-08.tex}
\subfile{Chapter-07/section-09.tex}
\subfile{Chapter-07/section-10.tex}

\end{document}
